\scanpage{552}
\mspage{463}

et que pour $\alpha<\beta$, on a
\[
M_\alpha\hookrightarrow M_\beta\qquad\text{injectif}
\]
Appliquons ceci au cas où $\alpha=\alpha_1$, donc on trouve $\operatorname{card} M_{\alpha_1}\geq\operatorname{card}\alpha_1=\pi_1$ d'autre part en \unsure{prenant|posant} $\beta=\alpha_1+1$, d'où $M_\beta\simeq M$, on trouve une injection
\[
M_{\alpha_1}\hookrightarrow M
\]
d'où $\pi_1=\operatorname{card}M_{\alpha_1}\leq\operatorname{card}M=\pi_0$, une contradiction, ce qui établit b$'$.

\deleted{\uline{\unsure{Corollaire}}} \deleted{\uline{\unsure{Scholie|Lemme}}} \inserted{Corollaire} 4.15.11. Soit $C$ une catégorie essentiellement petite, $\pi$ un cardinal
tel que $\pi>\operatorname{card}\,\unsure{\mathrm{red}|\mathrm{réd}}\,C$. \deleted{Alors \unsure{est|M}} \inserted{Soit} $\unsure{\mathcal M|M}\overset{\mathrm{def}}{=}\mathrm{Ind}_\pi C$ \deleted{$\unsure{=\mathcal M_\pi|\mathcal M_\pi}$}, alors $\mathcal M=\mathcal M_\pi$, \deleted{et} c'est une enveloppe de Karoubi de $C$.

\unsure{OPS|OPJ} $C$ réduite, donc $\operatorname{card} C<\pi$, et on applique le \deleted{\unsure{th}} corollaire précédent.

Ainsi, le \deleted{\unsure{th}} résultat litigieux 4.8.13 (p. 212) est finalement prouvé sous une forme légèrement plus faible, avec $\pi>\operatorname{card}\,\mathrm{red}\,C$ \unsure{ald|al d} $\pi\geq\operatorname{card}\,\mathrm{red}\,C$. Du $=$ coup, les applications de 4.8.13 qu'on a faites jusqu'à présent (notamment le lemme 4.15.3).
% archive sheet number: 551
